% Generated by tools/edn_latex.py from reports/edn.md.
% Do not edit: change reports/edn.md and run `make edn`.
\ednentry{
  id = {EDN-03},
  title = {New-market drift scenario: hold out AutoScout24 Spain instead of using DataMarket},
  date = {2026-09-22},
  milestone = {M1: Project Inception (data plan; shapes M6: Monitoring)},
  topic = {Dataset Selection, Monitoring},
  participants = {lukas2510, after supervisor feedback (S. del Rey)},
  decision = {Hold out all AutoScout24 \texttt{ES} listings (8,015, 6.8\,\%) from training, validation and calibration, and use them as the ``new market'' drift scenario for M6. The DataMarket Spanish second-hand car sample is dropped from the pipeline.},
  rationale = {\emph{Alternatives considered:}
\begin{itemize}[leftmargin=*, topsep=2pt]
\item \textbf{Option A (chosen): Hold out one AutoScout24 country (\texttt{ES}).}
\newline Pros: same schema, scrape date and portal as the training data, so detected drift comes from the market alone; every held-out listing has a price, so we can show real MAE and interval-coverage degradation, not only input drift; the same prices serve as delayed labels for the feedback loop; no mapping stage and no license issue; \texttt{ES} costs only 6.8\,\% of training rows (holding out \texttt{IT} would cost 20\,\%).
\newline Cons: less training data; the served model does not cover Spain until retrained; \texttt{country} becomes a category unseen in training, which the API must handle.
\item \textbf{Option B: Keep DataMarket as an optional stress test, restricted to makes both datasets share.}
\newline Pros: real external data, from a different portal; answers the supervisor's question with an actual experiment.
\newline Cons: market, time (2021 vs 2025), portal and brand mix still shift together, so a drift alarm cannot be attributed; needs a schema-mapping stage for \textasciitilde{}10 overlapping fields with Spanish labels and free-text colour; no license file and a commercial sample, so it cannot live in a public DVC remote.
\item \textbf{Option C: Drop DataMarket without a replacement.}
\newline Pros: simplest.
\newline Cons: leaves M6 without a realistic drift scenario.
\end{itemize}
\par
\emph{Why:} The supervisor asked us to check whether DataMarket is feasible as a new-market set given the feature mismatch. Our checks found that Spain is already in AutoScout24 (8,015 \texttt{ES} listings), 36.5\,\% of DataMarket rows are makes the model never sees and the shared mass-market makes have under 400 training rows each, and several factors shift at once (e.g. BMW median price 16,000 EUR in DataMarket vs 24,819 EUR in AutoScout24 \texttt{ES}). A drift alarm on DataMarket would therefore not be explainable. The M6 rubric asks for ``timely, actionable signals about data changes, model performance degradation'', and holding out \texttt{ES} gives a drift with a single known cause plus ground-truth prices to show the performance drop, at the cost of one filter in the split stage.},
  ai = {Alternative assessment, Recommendation},
  response = {Accepted},
  assessment = {The three options came from our own feasibility check. AI (Claude Code) explained the trade-offs against the M6 rubric and recommended Option A with \texttt{ES} as the held-out country (smaller training loss than \texttt{IT}), keeping \texttt{country} as a feature with unseen countries treated as unknown, and an optional synthetic drift scenario with a controlled cause. We accepted it because it answers the supervisor's concern with evidence, keeps the pipeline to one data source, and lets M6 show both input drift and performance degradation.},
  aievidence = {Claude Code session, 2026-09-22: prompt ``explain this to me again and give me a recommendation [...] what is feasible for our project and still in the scope of the course'' on project brief \S{}3.2; the response compared options A-C against the M6 rubric and recommended A with \texttt{ES}.},
  evidence = {\href{https://github.com/mlops-2627q1-mds-upc/MLOPS_Recommenditos/pull/13}{PR \#13}, project brief \S{}3.2 (\texttt{docs/\allowbreak{}docs/\allowbreak{}project-\allowbreak{}brief.\allowbreak{}md}).},
}
